\documentclass[11pt]{article}
 
% Latin Modern: the modern successor to Computer Modern
\usepackage[T1]{fontenc}
\usepackage{lmodern}
\usepackage[utf8]{inputenc}
 
\usepackage[margin=1in]{geometry}
\usepackage{parskip}      % space between paragraphs, no indent
\usepackage{enumitem}
\usepackage{setspace}
\singlespacing            % single-spaced body
 
\pagestyle{plain}
 
\begin{document}
 
% --- Header (top left) ---
\noindent
Daniel Kadyrov\\
\today\\
200-Hour Teacher Training 2026-2027\\
Sonic Yoga Center
 
\vspace{1em}
 
% --- Title: centered, body-text size, bold ---
\begin{center}
  \textbf{The Bhagavad Gita Reflection}
\end{center}
 
% --- Body ---
% There is 1 written reflection for The Bhagavad Gita Book

% It is 1 typewritten page (it can be up to 2 pages if necessary)

% Prompt: Write a 1-2 page reflection on The Bhagavad Gita Book. We encourage you to  cover at least one  theme from the book, it can focus on a few of the verses,  a whole chapter or the full text itself. Please contemplate the teachings and organize a reflection on the content. 

% Questions to consider: Why did you pick this theme? How does this concept apply to your daily life? To your practice? To your work? Does this compare/contrast with any other readings/teachings? What have you learned about Yoga and spirituality from this theme?  Why is this theme important in the history of Yoga. How does it relate to the other paths of Yoga.  In what ways can you teach and share this teaching to help others along the path? 

The Bhagavad Gita is beautiful epic 


I found the Bhagavad Gita to one be one of the most profound and remarkable books that I read this year and I will be ever thankful that the Sonic Yoga 200-Hour Teacher Training finally enticed me to read it as part of its curriculum.

I was deeply enthralled in the epic's 
 
\end{document}